\documentclass[12pt]{article}
\usepackage[utf8]{inputenc}
\usepackage[margin=0.75in]{geometry}
\usepackage{amsmath}
\usepackage{amssymb}
\usepackage{amsthm}
\usepackage{graphicx} % Required for inserting images
\usepackage{array}
\usepackage{tikz-cd}


\title{MAT1100 Homework 4}
\author{Aaron Avram}
\date{November 17 2026}

\begin{document}
\newtheorem{definition}{Definition}
\newtheorem{theorem}{Theorem}
\newtheorem{claim}{Claim}
\newtheorem{corollary}{Corollary}

\maketitle

\begin{align*}
    \textbf{Question 1}
\end{align*}

\begin{enumerate}
    \item[(i)]
    \begin{proof}
        $p(x) \in \ker \phi \iff p(0) = p(2) = 0$. This equivalent to saying
        that $x, x-2 \mid p(x)$ in $\mathbb Q[x]$ and it follows
        that $x(x-2) \mid p(x)$. Now suppose $q(x) \in \mathbb Q[x]$
        such that $p(x) = q(x)x(x-2)$, we wish to show $q(x) \in \mathbb Z[x]$. Without loss of generality $p$ is primitive,
        as if not prove the statement for $\frac{1}{g}p(x)$ where $g$ is the gcd of $p$'s
        coefficients and clearly this implies it for $p(x)$. Now there exists $m \in \mathbb Z$
        such that $mq(x)$ is an integer polynomial and therefore $a \in \mathbb Z$ such that
        $\frac{m}{a}q(x)$ is primitive. Then
        \[ \frac{m}{a}p(x) = (\frac{m}{a}q(x))(x)(x-2) \]
        all the polynomial's on the right hand side are integer polynomials and primitive, so their product is primitive.
        $p(x)$ is primitive, so the gcd of the coefficients of the product on the right hand side is $\frac{m}{a}$
        and this forces $\frac{m}{a} = 1$, i.e. $q(x) \in \mathbb Z[x]$ and thus $\ker \phi = \langle x(x-2) \rangle$.

        To compute the image consider $(a, b) \in \operatorname{im} \phi$ and suppose $\phi(p(x)) = (a, b)$
        then $p(0) = a$ and this implies that $p(x) = a_nx^n + \ldots a_1x + a$ for some $a_i \in \mathbb Z$.
        Thus
        \[ b = p(2) = a_n2^n + \ldots + a_12 + a = 2(a_n2^{n-1} + \ldots + a_1) + a \]
        and
        \[ b = a \mod 2 \]
        Therefore by the first isomorphism theorem
        \[ \operatorname{im} \phi = \{ (k, \ell) \in \mathbb Z^2 : k = \ell \mod 2 \} \cong R/\langle x(x-2) \rangle \]
        and the reverse inclusion holds as if $a = b \mod 2$ there exists $k$ such that $b - a = 2k$ and then $\phi(a + kx) = (a, b)$.
        I claim that $I + J = (x, 2)$. Clearly $x, 2 \in I + J$ as $2 = -(x-2) + x$
        and the reverse inclusion is trivial.

        Next $I \cap J = \langle x(x-2) \rangle$, as any polynomial in $I \cap J$ must be zero
        at $0$ and $2$ and by a similar argument to above this forces it to be divisible by $x(x-2)$.
        
        Finally, $IJ = \langle x(x-2) \rangle$ and this is clear as the product of two principal ideals
        is principal generated by the product of the generators.

        The map $R/IJ \to R/I \times R/J$ is not surjective as $(\overline{0}, \overline{1})$
        has no preimage in $R/IJ$ this is because if $f \in \mathbb Z[x]$ and modulo $I$, $f = 0$
        then $f(0) = 0$. However this implies that $f(2) = 0 \mod 2$. But $f(2) = 1 \mod J$
        so $f(x) = (x-2)p(x) + 1$ so $f(2) = 1$ which is a contradiction, thus no such $f$ exists.
    \end{proof}

    \item[(ii)]
    \begin{proof}
        $I + J$ is maximal as
        \[ R/(I + J) \cong \mathbb Z[x]/ \langle x, 2 \rangle \cong \mathbb Z/2\mathbb Z[x]/(x) \cong \mathbb Z/2\mathbb Z \]
        and $I, J \lneq I + J$ as $2 \notin I, J$ by degree considerations thus neither ideal can be maximal.
        Since $R/(I+J)$ is a field, it is in particular an integral domain therefore $R/I+J$ is also prime.

        Finally, $I, J$ are prime but not maximal as the evaluation maps $\operatorname{ev}_0$
        and $\operatorname{ev}_2$ give the isomorphisms
        \[ R/I \cong R/J \mathbb Z \]
        which is a domain but not a field. Where the fact that the kernels are exactly $IJ$
        follows from the same argument used to show $\ker \phi = IJ$.

        $I + J$ is not principal as if $\alpha \in \mathbb Z[x]$ such that $(\alpha) = (x, 2)$
        then $\alpha \mid x$ and $\alpha \mid 2$. Now the only elements that divide 2 are $\pm 1$ and $\pm 2$
        and clearly $\pm 2$ does not divide $x$ therefore $\alpha = \pm 1$. This implies that
        $I + J = (\pm 1) = \mathbb Z[x]$ which contradicts maximality.
    \end{proof}

    \item[(iii)]
    \begin{proof}
        By Gauss's lemma a non-constant polynomial in $\mathbb Z[x]$
        is irreducible if and only if it is primitive and irreducible in $\mathbb Q[x]$.
        Moreover, it is clear that any constant polynomial that is prime in $\mathbb Z$
        is irreducible in $\mathbb Z[x]$.

        Now if $f(x) \in \mathbb Z[x]$ then $g$ the gcd of $f$'s coefficients
        is unique up to units and $\frac{1}{g}f(x)$ is primitive. Now we may factor
        $\frac{1}{g}f(x)$ in $\mathbb Q[x]$ as a product of irreducible polynomials in $\mathbb Q[x]$
        $r_1(x), \ldots, r_k(x)$. For each $i$ there exists $m_i \in \mathbb Q$ such that $m_ir_i \in \mathbb Z[x]$
        and is primitive. Since this an associate of $r_i$ in $\mathbb Q[x]$ Gauss's lemma implies that
        $m_ir_i$ is irreducible in $\mathbb Z[x]$. Thus
        \[ \frac{1}{g}f = \prod_{i = 1}^k \frac{1}{m_i}(m_1r_1\cdot \ldots \cdot m_kr_k) \]
        now each $m_ir_i$ is primitive so their product is primitive, but we have that
        \[ \prod_{i = 1}^k(m_i)\Big(\frac{1}{g}f) = m_1r_1\cdot \ldots \cdot m_kr_k \]
        and therefore $\prod_{i = 1}^km_i = \pm 1$. Thus if $g = p_1\ldots p_\ell$ over $\mathbb Z$ then
        \[ f = (p_1\cdot \ldots p_\ell)(m_1r_1 \ldots m_kr_k) \]
        Now if $f$ has another factorization
        \[ f = \prod_{j}q_j \]
        then the constant factors $q_j$ multiply to get some integer $h$
        and since the non-constant factors must be primitive their product is primitive. This implies that
        the gcd of the coefficients of $f$ is equal to $g$ which is equal to $h$.
        Thus, by uniqueness of factorization in $\mathbb Z$, $g = \pm h$ and therefore
        $\frac{1}{g}f$ is a primitive polynomial with two factorizations into irreducible polynomials
        in $\mathbb Z[x]$. We can lift this factorization to $\mathbb Q[x]$ as primitive integer polynomials
        are irreducible over $\mathbb Q$. By uniqueness of factorization in $\mathbb Q[x]$
        this implies that the non-constant $q_j$ correspond exactly to the $m_ir_i$ up to units in $\mathbb Q[x]$
        but since all of these polynomials are primitive, if $q_j$ and $m_ir_i$ are associates in $\mathbb Q[x]$ they can only differ by $\pm 1$
        and so are also associates over $\mathbb Z$. Hence the factorization of $f$ is unique up to units.

        $\mathbb Z[x]$ is not a PID as by above $I + J$ is not principal.

        Finally it is clear that $1$ is a greatest common divisor of $x$ and $x - 2$
        as if $\alpha \mid x$ and $\alpha \mid x-2$ then $\alpha \mid 2$ which implies $\alpha = \pm 1$
        or $\alpha =  \pm 2$ but $\pm 2$ does not divide $x$ therefore $\alpha = \pm 1$
        and it is clear that $1$ divides both $x$ and $x-2$ and $-1$ is associate to $1$, therefore
        $1$ is a gcd of $x$ and $x-2$. 
        However there does not exist $A, B \in \mathbb Z[x]$
        such that
        \[ Ax + B(x - 2) = 1 \]
        as this would imply that $1 \in (x, x-2) = (x, 2)$ but $(x, 2)$ is maximal so this is not possible.
    \end{proof}
\end{enumerate}
\end{document}